%% kapa: exam=1 round=28 year=2017 period=1 form=A subject=economics subject_ko=경제학원론 seq=25 number=65
%% source: Q-Net 감정평가사 기출문제 (cst00309) artlSeq=5199603 file=extracted/1008/5199603/2017년 제28회 감정평가사 1차시험 문제지/1교시 A형.pdf
%% conversion: kordoc 4.24.51 -> tools/gen1.py (auto); review=proofread-sonnet; flags=none
\begin{question}{65}
모든 시장이 완전경쟁적인 甲국의 총생산함수는 $Y=AL^{\alpha}K^{1-\alpha}$이다. 甲국 경제에 관한 설명으로 옳은 것은? (단, $Y$는 총생산량, $L$은 노동투입량, $K$는 자본투입량, $A$는 총요소생산성이고, $A>0$, $0<\alpha<1$, 생산물 가격은 1이다.)
\begin{choices}
  \item 총생산량이 100이고 $\alpha =0.7$ 일 때, 자본에 귀속되는 자본소득은 70이다.
  \item $A$가 불변이고 $\alpha =0.7$ 일 때, 노동투입량이 3 \% 증가하고 자본투입량이 5 \% 증가하면 총생산량은 3 \% 증가한다.
  \item $A=3$일 때 노동과 자본의 투입량이 2 \%로 동일하게 증가하면 총생산량은 2 \%로 증가한다.
  \item 노동의 투입량이 5 \% 증가할 때 자본의 투입량도 5 \% 증가된다면, 노동의 한계 생산물은 변한다.
  \item $A$가 1 \% 증가하고 노동과 자본의 투입량이 모두 동일하게 2 \% 증가할 때, $\alpha$가 0.5보다 크다면 총생산량의 증가율은 5 \%이다.
\end{choices}
\end{question}
\answer{3}
